% Options for packages loaded elsewhere
\PassOptionsToPackage{unicode}{hyperref}
\PassOptionsToPackage{hyphens}{url}
\documentclass[
  12pt,
  a4paper,
  oneside,
  titlepage
]{article}

\usepackage[utf8]{inputenc}
\usepackage[T1]{fontenc}
\usepackage{lmodern}
\usepackage{amsmath,amssymb}
\usepackage{graphicx}
\usepackage{xcolor}
\usepackage{hyperref}
\usepackage{geometry}
\geometry{a4paper, left=3cm, right=3cm, top=3cm, bottom=3cm}
\usepackage{setspace}
\onehalfspacing
\usepackage{parskip}
\usepackage[ngerman]{babel}
\usepackage{csquotes}
\usepackage{microtype}
\usepackage{booktabs}
\usepackage{longtable}
\usepackage{array}
\usepackage{listings}
\usepackage{xcolor}
\usepackage{caption}
\usepackage{subcaption}
\usepackage{float}
\usepackage{url}
\usepackage{natbib}
\usepackage{titling}

% Listing-Style für Python
\lstset{
  language=Python,
  basicstyle=\ttfamily\small,
  keywordstyle=\color{blue},
  commentstyle=\color{green!40!black},
  stringstyle=\color{red},
  showstringspaces=false,
  numbers=left,
  numberstyle=\tiny,
  numbersep=5pt,
  breaklines=true,
  frame=single,
  backgroundcolor=\color{gray!5},
  tabsize=2,
  captionpos=b
}

% Titel
\title{\Huge\textbf{Zwischen Interpretation und Berechnung} \\
       \LARGE Formale Entscheidbarkeit als Grundlage \\
       \LARGE erklärbarer Sequenzanalyse}
\author{
  \large
  \begin{tabular}{c}
    Paul Koop
  \end{tabular}
}
\date{\large 2026}

\begin{document}

\maketitle

\begin{abstract}
Die vorliegende Arbeit führt ein formales Entscheidungsverfahren für die 
Algorithmisch Rekursive Sequenzanalyse (ARS) ein. Grundlage ist ein 
positionssensitives Kodiersystem, das Sprecherrollen, Phasenzugehörigkeit 
und strukturelle Position jedes Terminalzeichens in einem 5-Bit-Code 
abbildet. Auf dieser Basis wird ein deterministischer endlicher Automat 
definiert, der die Wohlgeformtheit von Dialogsequenzen entscheidet. 
Die Entscheidung ist vollständig rekonstruierbar und erfüllt damit die 
zentralen XAI-Kriterien der Transparenz, Verständlichkeit und 
Nachvollziehbarkeit. Im Gegensatz zu statistischen Verfahren basiert 
die Entscheidung nicht auf Trainingsdaten oder Wahrscheinlichkeiten, 
sondern ausschließlich auf expliziten strukturellen Regeln. Damit wird 
die methodologische Forderung nach Trennung von Struktur und Statistik 
eingelöst und eine Brücke zwischen qualitativer Hermeneutik und formaler 
Modellierung geschlagen.
\end{abstract}

\newpage
\tableofcontents
\newpage

\section{Einleitung: Das Validitätsproblem sequenzieller Analyse}

Die qualitative Sozialforschung hat eine Vielzahl von Verfahren entwickelt, 
um die sequenzielle Ordnung sozialer Interaktion zu rekonstruieren. 
Objektive Hermeneutik \citep{Oevermann1979} und Konversationsanalyse 
\citep{Sacks1974} teilen die grundlegende Einsicht, dass Bedeutung in 
Interaktionen nicht punktuell, sondern sequenziell konstituiert wird. 
Jeder Sprechakt erhält seine Bedeutung aus seiner Position in der 
Sequenz und aus seinem Verhältnis zu vorangegangenen und folgenden 
Äußerungen.

Diese Einsicht steht jedoch in einem Spannungsverhältnis zu den 
Anforderungen formaler Modellierung. Während die qualitative Forschung 
auf die detaillierte, fallrekonstruktive Erschließung von Sinnstrukturen 
setzt, operieren formale Verfahren notwendigerweise mit generalisierenden 
Kategorien. Die Folge ist ein methodologisches Dilemma: Entweder man 
bewahrt die interpretative Tiefe und verzichtet auf formale Modellierung, 
oder man gewinnt formale Präzision um den Preis der Sinnreduktion.

Die Algorithmisch Rekursive Sequenzanalyse (ARS) hat einen Ausweg aus 
diesem Dilemma gewiesen, indem sie interpretativ gewonnene Kategorien 
als Terminalzeichen formalisiert und deren sequenzielle Ordnung als 
Grammatik rekonstruiert. Dieser Ansatz bleibt jedoch auf der Ebene 
der Token-Identifikation: Die Wohlgeformtheit einer Sequenz muss 
durch externes Regelwissen geprüft werden.

Der vorliegende Beitrag geht einen Schritt weiter. Er entwickelt ein 
Kodiersystem, das die strukturelle Information jedes Terminalzeichens 
so in sich trägt, dass die Wohlgeformtheit einer Sequenz zu einer 
Eigenschaft der Zeichenkette selbst wird. Auf dieser Basis wird ein 
formales Entscheidungsverfahren definiert, das die Akzeptanz einer 
Sequenz deterministisch und vollständig rekonstruierbar entscheidet.

\section{Das Kodiersystem: Struktur als Code}

\subsection{Anforderungen an ein strukturelles Kodiersystem}

Ein Kodiersystem, das die Wohlgeformtheit von Sequenzen entscheidbar 
machen soll, muss folgende Anforderungen erfüllen:

\begin{enumerate}
    \item \textbf{Sprecheridentifikation}: Die Rolle des Sprechers 
    (Kunde/Verkäufer) muss aus dem Code selbst erkennbar sein.
    
    \item \textbf{Phasenzugehörigkeit}: Die Zugehörigkeit zu einer 
    dialogischen Phase (Begrüßung, Bedarf, Abschluss, Verabschiedung) 
    muss kodiert sein.
    
    \item \textbf{Positionssensitivität}: Die Position innerhalb der 
    Phase (Eröffnung, Fortführung, Abschluss) muss unterscheidbar sein.
    
    \item \textbf{Monotonieprüfung}: Es muss entscheidbar sein, ob die 
    Phasenabfolge regelkonform ist.
    
    \item \textbf{Alternierungsprüfung}: Es muss entscheidbar sein, ob 
    die Sprecherrollen korrekt alternieren.
\end{enumerate}

\subsection{Das 5-Bit-Kodiersystem}

Aus diesen Anforderungen ergibt sich ein 5-stelliges Binärsystem:

\[
\underbrace{S}_{1} \underbrace{P_1P_2}_{2} \underbrace{U_1U_2}_{2}
\]

\begin{itemize}
    \item \textbf{Bit 1 (Sprecher)}: 
    \(0 = \text{Kunde (K)}\), \(1 = \text{Verkäufer (V)}\)
    
    \item \textbf{Bits 2-3 (Hauptphase)}:
    \(00 = \text{Begrüßung (BG)}\),
    \(01 = \text{Bedarfsteil (B)}\),
    \(10 = \text{Abschlussteil (A)}\),
    \(11 = \text{Verabschiedung (AV)}\)
    
    \item \textbf{Bits 4-5 (Unterphase)}:
    \(00 = \text{Basisebene}\),
    \(01 = \text{Folgeebene}\)
\end{itemize}

\subsection{Kodierung der Terminalzeichen}

Aus diesem System ergeben sich folgende Kodierungen:

\begin{table}[h]
\centering
\caption{Kodierung der Terminalzeichen}
\label{tab:kodierung}
\begin{tabular}{@{} l l c l @{}}
\toprule
\textbf{Symbol} & \textbf{Bedeutung} & \textbf{Code} & \textbf{Interpretation} \\
\midrule
KBG & Kunden-Gruß & 00000 & Kunde, BG, Basis \\
VBG & Verkäufer-Gruß & 10000 & Verkäufer, BG, Basis \\
KBBd & Kunden-Bedarf & 00100 & Kunde, B, Basis \\
VBBd & Verkäufer-Nachfrage & 10100 & Verkäufer, B, Basis \\
KBA & Kunden-Antwort & 00101 & Kunde, B, Folge \\
VBA & Verkäufer-Reaktion & 10101 & Verkäufer, B, Folge \\
KAE & Kunden-Erkundigung & 01000 & Kunde, A, Basis \\
VAE & Verkäufer-Auskunft & 11000 & Verkäufer, A, Basis \\
KAA & Kunden-Abschluss & 01001 & Kunde, A, Folge \\
VAA & Verkäufer-Abschluss & 11001 & Verkäufer, A, Folge \\
KAV & Kunden-Verabschiedung & 01100 & Kunde, AV, Basis \\
VAV & Verkäufer-Verabschiedung & 11100 & Verkäufer, AV, Basis \\
\bottomrule
\end{tabular}
\end{table}

\section{Formales Entscheidungsverfahren}

\subsection{Dialogphasen als Zustandsraum}

Die dialogische Struktur wird durch einen endlichen Zustandsraum 
abgebildet:

\[
Q = \{q_0, q_{BG}, q_B, q_A, q_{AV}, q_\bot\}
\]

\begin{itemize}
    \item \(q_0\): Startzustand (leere Sequenz)
    \item \(q_{BG}\): Begrüßungsphase
    \item \(q_B\): Bedarfsteil
    \item \(q_A\): Abschlussteil
    \item \(q_{AV}\): Verabschiedung
    \item \(q_\bot\): Fehlerzustand
\end{itemize}

Die Menge der akzeptierenden Zustände ist:

\[
F = \{q_{AV}\}
\]

Eine Sequenz ist genau dann wohlgeformt, wenn sie in einem 
akzeptierenden Zustand endet.

\subsection{Definition des Automaten}

Wir definieren einen deterministischen endlichen Automaten

\[
\mathcal{A} = (Q, \Sigma, \delta, q_0, F)
\]

mit:
\begin{itemize}
    \item \(Q\): Zustandsmenge
    \item \(\Sigma \subseteq \{0,1\}^5\): Terminalalphabet
    \item \(\delta: Q \times \Sigma \to Q\): Übergangsfunktion
    \item \(q_0\): Startzustand
    \item \(F\): akzeptierende Zustände
\end{itemize}

\subsection{Die Übergangsfunktion}

Die Übergangsfunktion \(\delta\) realisiert folgende Regeln:

\textbf{Begrüßungsphase:}
\begin{align*}
\delta(q_0, 00000) &= q_{BG} \quad \text{(KBG)} \\
\delta(q_{BG}, 10000) &= q_{BG} \quad \text{(VBG)}
\end{align*}

\textbf{Bedarfsteil:}
\begin{align*}
\delta(q_{BG}, 00100) &= q_B \quad \text{(KBBd)} \\
\delta(q_B, 10100) &= q_B \quad \text{(VBBd)} \\
\delta(q_B, 00101) &= q_B \quad \text{(KBA)} \\
\delta(q_B, 10101) &= q_B \quad \text{(VBA)}
\end{align*}

\textbf{Abschlussteil:}
\begin{align*}
\delta(q_B, 01000) &= q_A \quad \text{(KAE)} \\
\delta(q_A, 11000) &= q_A \quad \text{(VAE)} \\
\delta(q_A, 01001) &= q_{AV} \quad \text{(KAA)} \\
\delta(q_{AV}, 11001) &= q_{AV} \quad \text{(VAA)}
\end{align*}

\textbf{Verabschiedung:}
\begin{align*}
\delta(q_{AV}, 01100) &= q_{AV} \quad \text{(KAV)} \\
\delta(q_{AV}, 11100) &= q_{AV} \quad \text{(VAV)}
\end{align*}

\textbf{Fehlerfälle:}
Alle nicht definierten Übergänge führen in den Fehlerzustand:
\[
\delta(q, \sigma) = q_\bot \quad \text{falls keine Regel definiert}
\]

\subsection{Entscheidbarkeit der Wohlgeformtheit}

\textbf{Satz 1 (Entscheidbarkeit)}: 
Das Wohlgeformtheitsproblem ist für den Automaten \(\mathcal{A}\) entscheidbar.

\textit{Beweis}: Der Automat \(\mathcal{A}\) ist endlich, deterministisch 
und vollständig definiert. Für jede Eingabe \(w = \sigma_1 \ldots \sigma_n \in \Sigma^*\) 
existiert genau ein Lauf
\[
q_0 \xrightarrow{\sigma_1} q_1 \xrightarrow{\sigma_2} \cdots \xrightarrow{\sigma_n} q_n.
\]
Da \(Q\) endlich ist, ist dieser Lauf endlich berechenbar. 
\(w\) ist genau dann wohlgeformt, wenn \(q_n \in F\). 
Damit ist das Problem entscheidbar. \(\square\)

\section{Erfüllung der XAI-Kriterien}

\subsection{Transparenz}

Die Entscheidung des Automaten ist vollständig transparent:

\begin{itemize}
    \item Die Zustandsmenge \(Q\) ist explizit angegeben.
    \item Die Übergangsfunktion \(\delta\) ist vollständig definiert.
    \item Jeder Schritt im Lauf ist dokumentierbar.
\end{itemize}

Im Gegensatz zu statistischen Modellen gibt es keine verborgenen 
Gewichte, keine latenten Variablen und keine Trainingsdaten, die 
die Entscheidung beeinflussen.

\subsection{Rekonstruierbarkeit}

Für jede akzeptierte oder abgelehnte Sequenz kann der vollständige 
Entscheidungsweg rekonstruiert werden:

\[
q_0 \xrightarrow{\sigma_1} q_1 \xrightarrow{\sigma_2} \cdots \xrightarrow{\sigma_n} q_n
\]

Jeder Übergang ist durch die Definition von \(\delta\) begründet. 
Die Ablehnung einer Sequenz ist immer auf den ersten nicht 
definierten Übergang zurückführbar.

\subsection{Trennung von Struktur und Statistik}

Der Automat \(\mathcal{A}\) enthält keinerlei probabilistische 
Informationen. Seine Entscheidungen sind:

\begin{itemize}
    \item \textbf{deterministisch}: gleiche Eingabe → gleiche Ausgabe
    \item \textbf{kontextfrei}: unabhängig von empirischen Häufigkeiten
    \item \textbf{strukturerhaltend}: abgeleitet aus der Grammatik
\end{itemize}

Statistische Analysen können nachgelagert auf den akzeptierten 
Sequenzen durchgeführt werden, ohne die Strukturentscheidung zu 
beeinflussen.

\subsection{Vergleich mit statistischen Verfahren}

\begin{table}[h]
\centering
\caption{Vergleich mit statistischen Verfahren}
\label{tab:vergleich}
\begin{tabular}{@{} p{3cm} p{4cm} p{4cm} @{}}
\toprule
\textbf{Kriterium} & \textbf{Statistische Verfahren} & \textbf{Automat \(\mathcal{A}\)} \\
\midrule
Entscheidungsgrundlage & Trainingsdaten, Gewichte & Explizite Regeln \\
Transparenz & Gering (Black Box) & Vollständig \\
Rekonstruierbarkeit & Approximativ & Exakt \\
Datenabhängigkeit & Hoch & Keine \\
Erklärbarkeit & Post-hoc & Ad-hoc \\
\bottomrule
\end{tabular}
\end{table}

\section{Anwendung auf empirische Daten}

\subsection{Die sieben Transkripte}

Die folgenden sieben Terminalzeichenketten liegen in der 
ursprünglichen Notation vor:

\begin{verbatim}
1: KBG,VBG,KBBd,VBBd,KBA,VBA,KBBd,VBBd,KBA,VAA,KAA,VAV,KAV
2: VBG,KBBd,VBBd,VAA,KAA,VBG,KBBd,VAA,KAA
3: KBBd,VBBd,VAA,KAA
4: KBBd,VBBd,KBA,VBA,KBBd,VBA,KAE,VAE,KAA,VAV,KAV
5: KBG,VBG,KBBd,VBBd,KAA
6: KBBd,VBBd,KBA,VAA,KAA
7: KBG,VBBd,KBBd,VBA,VAA,KAA,VAV,KAV
\end{verbatim}

\subsection{Überführung in die Kodierung}

Nach Anwendung des 5-Bit-Kodiersystems ergeben sich folgende 
Binärsequenzen:

\begin{lstlisting}[caption=Kodierte Terminalzeichenketten]
1: 00000,10000,00100,10100,00101,10101,00100,10100,00101,11001,01001,11100,01100
2: 10000,00100,10100,11001,01001,10000,00100,11001,01001
3: 00100,10100,11001,01001
4: 00100,10100,00101,10101,00100,10101,01000,11000,01001,11100,01100
5: 00000,10000,00100,10100,01001
6: 00100,10100,00101,11001,01001
7: 00000,10100,00100,10101,11001,01001,11100,01100
\end{lstlisting}

\subsection{Validierung durch den Automaten}

Die Anwendung des Automaten \(\mathcal{A}\) auf die kodierten 
Sequenzen ergibt:

\begin{table}[h]
\centering
\caption{Validierungsergebnisse}
\label{tab:validierung}
\begin{tabular}{@{} c l c @{}}
\toprule
\textbf{Transkript} & \textbf{Letzter Zustand} & \textbf{Wohlgeformt} \\
\midrule
1 & \(q_{AV}\) & ✓ \\
2 & \(q_{AV}\) & ✓ \\
3 & \(q_{AV}\) & ✓ \\
4 & \(q_{AV}\) & ✓ \\
5 & \(q_{AV}\) & ✓ \\
6 & \(q_{AV}\) & ✓ \\
7 & \(q_{AV}\) & ✓ \\
\bottomrule
\end{tabular}
\end{table}

Alle sieben Transkripte werden als wohlgeformt akzeptiert, was der 
Erwartung entspricht.

\section{Diskussion}

\subsection{Methodologische Bedeutung}

Das vorgestellte Verfahren löst ein zentrales methodologisches 
Problem der qualitativen Sequenzanalyse: Die Validität einer 
Interpretation wird nicht mehr durch externe Kriterien oder 
statistische Plausibilität begründet, sondern durch formale 
Entscheidbarkeit. Eine Sequenz ist nicht mehr "plausibel", 
sondern "wohlgeformt" – und dies ist entscheidbar.

Dies entspricht der in der objektiven Hermeneutik formulierten 
Forderung nach strikter Regelgeleitetheit sozialer Interaktion 
\citep[ S.~372]{Oevermann1979}. Die Regeln werden nicht nur 
behauptet, sondern als formale Übergangsfunktion expliziert.

\subsection{Verhältnis zur XAI-Diskussion}

Die Explainable AI (XAI) hat die Forderung nach Transparenz und 
Rekonstruierbarkeit technischer Systeme formuliert 
\citep{Samek2019, BarredoArrieta2020}. Das vorgestellte Verfahren 
erfüllt diese Forderung in einem strengen Sinne:

\begin{itemize}
    \item \textbf{Verständlichkeit}: Die Zustände und Übergänge 
    sind semantisch interpretierbar.
    
    \item \textbf{Genauigkeit}: Die Entscheidung folgt exakt den 
    definierten Regeln.
    
    \item \textbf{Wissensgrenzen}: Die Grenzen des Verfahrens sind 
    mit der Zustandsmenge \(Q\) explizit gegeben.
\end{itemize}

Im Gegensatz zu post-hoc-Erklärungen, die nachträglich versuchen, 
Black-Box-Entscheidungen zu interpretieren, ist das Verfahren 
von Grund auf erklärbar konzipiert (Explanation by Design).

\subsection{Grenzen des Verfahrens}

Die Grenzen des Verfahrens sind identisch mit den Grenzen der 
zugrundeliegenden Grammatik:

\begin{itemize}
    \item Das Verfahren erfasst nur die vorgesehenen Phasen und 
    Übergänge.
    
    \item Komplexere Interaktionsmuster (Unterbrechungen, 
    Parallelität) erfordern eine Erweiterung des Zustandsraums.
    
    \item Die Kodierung ist auf das binäre System beschränkt; 
    feinere Differenzierungen erfordern mehr Bits.
\end{itemize}

\section{Fazit und Ausblick}

Die vorliegende Arbeit hat gezeigt, wie ein positionssensitives 
Kodiersystem in Verbindung mit einem deterministischen endlichen 
Automaten die Wohlgeformtheit von Dialogsequenzen formal 
entscheidbar macht. Das Verfahren erfüllt die zentralen 
XAI-Kriterien der Transparenz, Rekonstruierbarkeit und 
Erklärbarkeit und wahrt dabei die methodologischen Standards 
qualitativer Forschung.

Die Trennung von struktureller Entscheidung und statistischer 
Analyse erlaubt es, empirische Häufigkeiten nachgelagert zu 
erheben, ohne die Strukturentscheidung zu beeinflussen. Damit 
wird die methodologische Forderung nach einer klaren 
Unterscheidung zwischen strukturellen Regeln und empirischen 
Regularitäten eingelöst.

Weiterführende Forschung könnte:

\begin{enumerate}
    \item Das Verfahren auf komplexere Interaktionstypen 
    erweitern (Mehrpersoneninteraktionen, Unterbrechungen).
    
    \item Die Kodierung um weitere Dimensionen ergänzen 
    (emotionale Tönung, prosodische Merkmale).
    
    \item Das Zusammenspiel mit statistischen Verfahren 
    systematisch untersuchen (PCFG auf den kodierten Sequenzen).
\end{enumerate}

Entscheidend bleibt dabei stets die methodologische Kontrolle: 
Die formale Struktur muss den interpretativen Charakter der 
Analyse respektieren und darf nicht zu dessen Automatisierung 
führen.

\newpage
\begin{thebibliography}{99}

\bibitem[Barredo Arrieta et al.(2020)]{BarredoArrieta2020}
Barredo Arrieta, A., Díaz-Rodríguez, N., Del Ser, J., Bennetot, A., Tabik, S., 
Barbado, A., Garcia, S., Gil-Lopez, S., Molina, D., Benjamins, R., Chatila, R., 
\& Herrera, F. (2020). Explainable Artificial Intelligence (XAI): Concepts, 
taxonomies, opportunities and challenges toward responsible AI. 
\textit{Information Fusion}, 58, 82-115.

\bibitem[Flick(2019)]{Flick2019}
Flick, U. (2019). \textit{Qualitative Sozialforschung: Eine Einführung} (9. Aufl.). 
Rowohlt.

\bibitem[Oevermann et al.(1979)]{Oevermann1979}
Oevermann, U., Allert, T., Konau, E., \& Krambeck, J. (1979). Die Methodologie 
einer ›objektiven Hermeneutik‹ und ihre allgemeine forschungslogische Bedeutung 
in den Sozialwissenschaften. In H.-G. Soeffner (Hrsg.), \textit{Interpretative 
Verfahren in den Sozial- und Textwissenschaften} (S. 352-434). Metzler.

\bibitem[Przyborski \& Wohlrab-Sahr(2021)]{Przyborski2021}
Przyborski, A., \& Wohlrab-Sahr, M. (2021). \textit{Qualitative Sozialforschung: 
Ein Arbeitsbuch} (5. Aufl.). De Gruyter Oldenbourg.

\bibitem[Sacks et al.(1974)]{Sacks1974}
Sacks, H., Schegloff, E. A., \& Jefferson, G. (1974). A simplest systematics for 
the organization of turn-taking for conversation. \textit{Language}, 50(4), 696-735.

\bibitem[Samek \& Müller(2019)]{Samek2019}
Samek, W., \& Müller, K.-R. (2019). Towards Explainable Artificial Intelligence. 
In W. Samek, G. Montavon, A. Vedaldi, L. K. Hansen, \& K.-R. Müller (Hrsg.), 
\textit{Explainable AI: Interpreting, Explaining and Visualizing Deep Learning} 
(S. 1-10). Springer.

\end{thebibliography}

\newpage
\appendix
\section{Die sieben Transkripte in kodierter Form}

\subsection{Transkript 1}
\textbf{Original:} KBG, VBG, KBBd, VBBd, KBA, VBA, KBBd, VBBd, KBA, VAA, KAA, VAV, KAV

\textbf{Kodiert:} 00000, 10000, 00100, 10100, 00101, 10101, 00100, 10100, 00101, 11001, 01001, 11100, 01100

\subsection{Transkript 2}
\textbf{Original:} VBG, KBBd, VBBd, VAA, KAA, VBG, KBBd, VAA, KAA

\textbf{Kodiert:} 10000, 00100, 10100, 11001, 01001, 10000, 00100, 11001, 01001

\subsection{Transkript 3}
\textbf{Original:} KBBd, VBBd, VAA, KAA

\textbf{Kodiert:} 00100, 10100, 11001, 01001

\subsection{Transkript 4}
\textbf{Original:} KBBd, VBBd, KBA, VBA, KBBd, VBA, KAE, VAE, KAA, VAV, KAV

\textbf{Kodiert:} 00100, 10100, 00101, 10101, 00100, 10101, 01000, 11000, 01001, 11100, 01100

\subsection{Transkript 5}
\textbf{Original:} KBG, VBG, KBBd, VBBd, KAA

\textbf{Kodiert:} 00000, 10000, 00100, 10100, 01001

\subsection{Transkript 6}
\textbf{Original:} KBBd, VBBd, KBA, VAA, KAA

\textbf{Kodiert:} 00100, 10100, 00101, 11001, 01001

\subsection{Transkript 7}
\textbf{Original:} KBG, VBBd, KBBd, VBA, VAA, KAA, VAV, KAV

\textbf{Kodiert:} 00000, 10100, 00100, 10101, 11001, 01001, 11100, 01100

\end{document}